\subsection{Generalised Context Model}

We modelled loanword category assignment using the Generalised Context Model \citep[GCM;][]{nosofsky1988}. The GCM is an exemplar model: rather than representing a category by a prototype, it stores individual category members and assigns novel items to categories in proportion to their summed similarity to each category's exemplars.

\paragraph{Exemplar set.}
The training set comprised \(N = 106\) real Hungarian loanwords with known harmonic class. Words of Latin and French origin were assigned to category \textsc{front}; words of German and Yiddish origin were assigned to category \textsc{back}.

\paragraph{Phonological distance.}
\textbf{[TO FILL IN]}

\paragraph{Pairwise similarity.}
Similarity between a test item \(w\) and training exemplar \(e_j\) was computed as

\begin{equation}
  s_{wj} = \exp\!\left(-\frac{d_{wj}}{s}\right)^{\!p}
  \label{eq:gcm-sim}
\end{equation}

\noindent where \(d_{wj}\) is the phonological distance between \(w\) and \(e_j\), \(s > 0\) is a sensitivity parameter controlling the rate of similarity decay with distance, and \(p \in \{1, 2\}\) is the Minkowski exponent (\(p = 1\): city-block; \(p = 2\): Euclidean).

\paragraph{Category probability.}
The GCM prediction for a test item \(w\) was the probability of category \textsc{back},

\begin{equation}
  P(\textsc{back} \mid w) = \frac{\displaystyle\sum_{j \in \textsc{back}} s_{wj}}{\displaystyle\sum_{j=1}^{N} s_{wj}}
  \label{eq:gcm-prob}
\end{equation}

Because the two categories are exhaustive and complementary, \(P(\textsc{front} \mid w) = 1 - P(\textsc{back} \mid w)\).

\paragraph{Hyperparameter search.}
We conducted a grid search over \(s \in \{0.10, 0.35, \ldots, 4.85\}\) (step~0.25; 20~values) and \(p \in \{1, 2\}\), yielding 40 parameter combinations. For each combination, \(P(\textsc{back} \mid w)\) was computed for every test nonword and merged with the trial-level acceptance data. A mixed-effects logistic regression was fitted with \(P(\textsc{back})\) as the sole fixed-effect predictor and random intercepts by participant and by item:

\begin{equation}
  \operatorname{logit} P(\mathrm{accept}_{ij})
    = \beta_0 + \beta_1\,P(\textsc{back} \mid w_j) + u_i + v_j,
  \qquad
  u_i \sim \mathcal{N}(0,\,\sigma_u^2),\quad v_j \sim \mathcal{N}(0,\,\sigma_v^2)
  \label{eq:gcm-search}
\end{equation}

Root mean squared error (RMSE) between fitted acceptance probabilities and observed binary responses was computed for each combination; the pair minimising RMSE was selected.

\paragraph{Final model.}
Using the best-fitting hyperparameters, \(P(\textsc{back})\) was recomputed for all test items and $z$-scored prior to modelling. The final model was estimated in a Bayesian framework using \texttt{brms} \citep{burkner2017}, with a Bernoulli likelihood, default priors, and four MCMC chains. The primary inferential statistic was the posterior probability that $\beta_1 > 0$, i.e.\ that greater \textsc{back}-category similarity increases acceptance.
